\documentclass[11pt,a4paper] {article}
\usepackage[utf8] {inputenc}
\usepackage[polish] {babel}
\usepackage[T1] {fontenc}
\usepackage{graphicx}
\usepackage{hyperref}
\usepackage{listings}
\usepackage{xcolor}
\usepackage{geometry}
\usepackage{float}
\usepackage{amsmath}
\usepackage{longtable}
\usepackage{caption}

\geometry{margin=2.5cm}

\lstset{
  basicstyle=\ttfamily\small,
  breaklines=false,
  columns=fullflexible,
  frame=single,
  numbers=left,
  numberstyle=\tiny,
  language=json
}

\title{Module \& Interface Reference}
\author{Fire Simulation System}
% \date{\today}

\begin{document}

\maketitle
% \tableofcontents
\newpage

\section{Wprowadzenie}

\subsection{Cel dokumentu}

Dokument zawiera najbardziej techniczne szczegóły dotyczące API, kontraktów RabbitMQ, formatów JSON oraz interfejsów między modułami. 

\subsection{Zakres dokumentu}

\begin{itemize}
    \item REST API wszystkich serwisów
    \item Kontrakty RabbitMQ (kolejki, routing keys, formaty wiadomości)
    \item Format JSON dla konfiguracji i danych
    \item Wersjonowanie i kompatybilność wsteczną
    \item Walidacja i obsługa błędów
\end{itemize}

\section{REST API}

\subsection{FFBackend API}

\begin{lstlisting}[
  float,
  placement=H,
  caption={Żądanie uruchomienia symulacji}
]
POST /simulation/run-simulation?interval=5
Content-Type: application/json

{
  "forestName": "Test Forest",
  "sectors": [
    {
      "sectorId": 1,
      "sectorType": "FOREST",
      "initialState": {
        "fireLevel": 0,
        "burnLevel": 0,
        "extinguishLevel": 0,
        "fireState": "INACTIVE"
      },
      "contours": [[20.5, 49.9], [20.6, 49.9]],
      "assignedBrigades": []
    }
  ],
  "fireBrigades": [
    {
      "fireBrigadeId": 1,
      "currentLocation": {"longitude": 20.5, "latitude": 49.9},
      "state": "IDLE"
    }
  ],
  "foresterPatrols": [
    {
      "foresterPatrolId": 1,
      "currentLocation": {"longitude": 20.6, "latitude": 49.8},
      "state": "PATROLLING"
    }
  ]
}
\end{lstlisting}

\textbf{Przykład odpowiedzi SSE:}
\begin{lstlisting}[
  float,
  placement=H,
  caption={Request Body}
]
data: 
{  
    "forestName":"Test Forest",
    "timestamp":"2026-01-14T17:30:53.000Z",
    "tick":0,
    "sectors":[...],
    "fireBrigades":[...],
    "foresterPatrols":[...],
    "recommendedActions":[...]
}

data: 
{
    "forestName":"Test Forest",
    "timestamp":"2026-01-14T17:30:53.500Z",
    "tick":1,...
}
\end{lstlisting}

\paragraph{POST /simulation/send-simulation-request}

Wysyła żądanie uruchomienia symulacji do serwisu symulacji.

\textbf{Request Body:} \texttt{Configuration} (JSON)

\textbf{Response:}
\begin{itemize}
    \item \texttt{200 OK}: "Configuration sent to simulation!"
    \item \texttt{503 Service Unavailable}: "Simulation service unreachable"
\end{itemize}

\paragraph{POST /simulation/orderFireBrigade}

Wysyła komendę do brygady pożarniczej.

\textbf{Request Body:}
\begin{lstlisting}[
  float,
  placement=H,
  caption={}
]
{
  "id": 1,
  "location": {
    "longitude": 20.5,
    "latitude": 49.9
  },
  "goToBase": false
}
\end{lstlisting}

\textbf{Response:}
\begin{itemize}
    \item \texttt{200 OK}: "Order forwarded to simulation"
    \item \texttt{202 Accepted}: "Simulation unreachable; order queued"
    \item \texttt{503 Service Unavailable}: "Simulation unreachable and enqueue failed"
\end{itemize}

\paragraph{POST /simulation/orderForestPatrol}

Wysyła komendę do patrolu leśnego. Format identyczny jak \texttt{orderFireBrigade}.

\paragraph{POST /simulation/assignBrigades}

Przypisuje wiele brygad do sektora.

\textbf{Request Body:}
\begin{lstlisting}[
  float,
  placement=H,
  caption={}
]
{
  "sectorId": 5,
  "assignedBrigades": [1, 2, 3]
}
\end{lstlisting}

\textbf{Response:}
\begin{itemize}
    \item \texttt{200 OK}: "AssignBrigades forwarded to simulation"
    \item \texttt{503 Service Unavailable}: "Simulation service unreachable"
\end{itemize}

\paragraph{POST /simulation/stop-simulation}

Zatrzymuje uruchomioną symulację.

\textbf{Request Body:} Puste lub \texttt{""}

\textbf{Response:}
\begin{itemize}
    \item \texttt{200 OK}: "Stop request sent!"
    \item \texttt{503 Service Unavailable}: "Simulation service unreachable"
\end{itemize}

\textbf{Nie w pełni zaimplementowane! Zobacz dokument issues}

\paragraph{POST /simulation/set-speed}

Ustawia prędkość symulacji.

\textbf{Query Parameters:}
\begin{itemize}
    \item \texttt{tickInterval} (Double, wymagany): Interwał ticka symulacji w sekundach (musi być > 0)
\end{itemize}

\textbf{Response:}
\begin{itemize}
    \item \texttt{200 OK}: "Simulation speed updated"
    \item \texttt{400 Bad Request}: "tickInterval must be > 0"
    \item \texttt{503 Service Unavailable}: "Simulation service unreachable"
\end{itemize}

\textbf{Nie w pełni zaimplementowane! Zobacz dokument issues}

\subsubsection{Legacy Endpoints}

Dla kompatybilności wstecznej dostępne są również:
\begin{itemize}
    \item \texttt{POST /orderFireBrigade} (alias dla \texttt{/simulation/orderFireBrigade})
    \item \texttt{POST /orderForestPatrol} (alias dla \texttt{/simulation/orderForestPatrol})
\end{itemize}

\subsubsection{Error Responses}

Standardowy format błędu:
\begin{lstlisting}[
  float,
  placement=H,
  caption={}
]
{
  "timestamp": "2026-01-14T17:30:53.000Z",
  "status": 503,
  "error": "Service Unavailable",
  "message": "Simulation service unreachable",
  "path": "/simulation/orderFireBrigade"
}
\end{lstlisting}

\subsubsection{CORS}

Wszystkie endpointy wspierają CORS z \texttt{Access-Control-Allow-Origin: *} (konfigurowalne przez \texttt{WebConfig}).

\subsection{FFConf API}

\texttt{http://localhost:31415} (konfigurowalne przez \texttt{CONFIG\_PORT})

\subsubsection{Endpoints}

\begin{itemize}
  \item \textbf{GET /api/v1/nodes} -- Pobiera listę węzłów konfiguracji.
  \item \textbf{GET /api/v1/nodes/\{nodeId\}} -- Pobiera szczegóły węzła konfiguracji.
  \item \textbf{POST /api/v1/nodes} -- Tworzy nowy węzeł konfiguracji.
  \item \textbf{GET /health} -- Sprawdza status serwisu.
  \item \textbf{POST /admin/reload-configs} - Przeładowuje konfiguracje z dysku
\end{itemize}

\textbf{Query Parameters:}
\begin{itemize}
    \item \texttt{force} (bool, opcjonalny): Jeśli \texttt{true}, czyści kolekcję przed importem
\end{itemize}

\section{RabbitMQ Contracts}

\subsection{Exchange Configuration}

Wszystkie kolejki używają exchange typu \texttt{topic} o nazwie \texttt{fire-simulation-exchange} (konfigurowalne przez \texttt{RABBITMQ\_EXCHANGE}).

\subsection{Aktywne kolejki (wg backendu)}

Poniższa lista pochodzi z \texttt{DeclaredQueues} w backendzie i obejmuje tylko realnie używane kolejki; pomijamy kolejki oznaczone w kodzie jako \textit{NOT USED} (np. LLM chat).

\subsubsection{Telemetry FROM simulator}
\begin{itemize}
  \item \texttt{simulation\_telemetry\_agents\_fire\_brigade}
  \item \texttt{simulation\_telemetry\_agents\_fire\_brigade\_batch}
  \item \texttt{simulation\_telemetry\_agents\_forester}
  \item \texttt{simulation\_telemetry\_agents\_forester\_batch}
  \item \texttt{simulation\_telemetry\_sensors\_temp\_humidity}
  \item \texttt{simulation\_telemetry\_sensors\_wind\_speed}
  \item \texttt{simulation\_telemetry\_sensors\_wind\_direction}
  \item \texttt{simulation\_telemetry\_sensors\_litter\_moisture}
  \item \texttt{simulation\_telemetry\_sensors\_co2}
  \item \texttt{simulation\_telemetry\_sensors\_pm2\_5}
  \item \texttt{simulation\_telemetry\_sensors\_camera}
  \item \texttt{simulation\_telemetry\_map\_sector\_state}
  \item \texttt{simulation\_telemetry\_map\_sector\_state\_fast}
\end{itemize}

\subsubsection{Support / rekomendacje / agregacja}
\begin{itemize}
  \item \texttt{support\_analysis\_requests}
  \item \texttt{support\_recommendations}
  \item \texttt{support\_data\_aggregated}
  \item \texttt{simulation\_recommendations}
  \item \texttt{support\_analytics\_insights}
\end{itemize}

\subsubsection{Kolejki nieużywane (w kodzie oznaczone jako NOT USED)}
\begin{itemize}
  \item \texttt{backend\_llm\_requests}, \texttt{backend\_llm\_responses}
  \item \texttt{support\_llm\_propositions}, \texttt{support\_llm\_propositions\_responses}
  \item \texttt{support\_llm\_requests}, \texttt{support\_llm\_responses}
  \item \texttt{backend\_tasks\_queue}, \texttt{backend\_data\_aggregated}, \texttt{backend\_commands\_user}
\end{itemize}

\textbf{Uwaga:} Nazwa exchange jest konfigurowalna (\texttt{RABBITMQ\_EXCHANGE}). W środowisku lokalnym w plikach \texttt{.local.env} może występować \texttt{fire\_updates};

\subsection{Kolejki przychodzące (FROM Simulation)}

\subsubsection{simulation\_telemetry\_agents\_fire\_brigade}

\textbf{Routing Key:} \texttt{simulation.telemetry.agents.fire\_brigade}

\textbf{Event Model:} \texttt{EvFireBrigade}

\textbf{Contract:}
\begin{lstlisting}[
  float,
  placement=H,
  caption={}
]
{
  "fireBrigadeId": 1,
  "sectorId": 5,
  "location": {
    "longitude": 20.5,
    "latitude": 49.9
  },
  "state": "EXTINGUISHING"
}
\end{lstlisting}

\textbf{Required Fields:}
\begin{itemize}
    \item \texttt{fireBrigadeId} (int): Unikalny identyfikator brygady
    \item \texttt{location} (Location): Aktualne współrzędne GPS
    \item \texttt{state} (FireBrigadeState): Aktualny stan (EXTINGUISHING, IDLE, MOVING, etc.)
\end{itemize}

\textbf{Optional Fields:}
\begin{itemize}
    \item \texttt{sectorId} (Integer): ID sektora jeśli znane, w przeciwnym razie rozwiązywane z location
\end{itemize}

\subsubsection{simulation\_telemetry\_agents\_forester}

\textbf{Routing Key:} \texttt{simulation.telemetry.agents.forester}

\textbf{Event Model:} \texttt{EvForestPatrol}

\textbf{Contract:}
\begin{lstlisting}[
  float,
  placement=H,
  caption={}
]
{
  "foresterPatrolId": 1,
  "sectorId": 3,
  "location": {
    "longitude": 20.6,
    "latitude": 49.8
  },
  "state": "PATROLLING"
}
\end{lstlisting}

\subsubsection{simulation\_telemetry\_map\_sector\_state}

\textbf{Routing Key:} \texttt{simulation.telemetry.map.sector\_state}

\textbf{Event Model:} \texttt{EvSectorState}

\textbf{Contract:}
\begin{lstlisting}[
  float,
  placement=H,
  caption={}
]
{
  "sectorId": 1,
  "fireLevel": 25.5,
  "burnLevel": 10.2,
  "extinguishLevel": 5.0
}
\end{lstlisting}

\textbf{Required Fields:}
\begin{itemize}
    \item \texttt{sectorId} (int): Identyfikator sektora
    \item \texttt{fireLevel} (double): Aktualna intensywność ognia (0-100)
    \item \texttt{burnLevel} (double): Procent spalenia (0-100)
    \item \texttt{extinguishLevel} (double): Postęp gaszenia (0-100)
\end{itemize}

\textbf{Optimization Note:} Powinny być wysyłane tylko zmienione sektory, nie wszystkie sektory za każdym razem.

\subsubsection{Czujniki}

Wszystkie kolejki czujników używają tego samego wzorca z routing opartym na lokalizacji.

\paragraph{simulation\_telemetry\_sensors\_temp\_humidity}

\textbf{Routing Key:} \texttt{simulation.telemetry.sensors.temp\_humidity}

\textbf{Event Model:} \texttt{EvTempAndAirHumiditySensor}

\textbf{Contract:}
\begin{lstlisting}[
  float,
  placement=H,
  caption={}
] 
{
  "sensorId": 1,
  "timestamp": "2026-01-14T17:30:53",
  "sensorType": "TEMP_AND_AIR_HUMIDITY",
  "location": {
    "longitude": 20.5,
    "latitude": 49.9
  },
  "data": {
    "temperature": 25.5,
    "airHumidity": 60.0
  }
}
\end{lstlisting}

\paragraph{simulation\_telemetry\_sensors\_wind\_speed}

\textbf{Routing Key:} \texttt{simulation.telemetry.sensors.wind\_speed}

\textbf{Event Model:} \texttt{EvWindSpeedSensor}

\textbf{Contract:}
\begin{lstlisting}[
  float,
  placement=H,
  caption={}
]
{
  "sensorId": 2,
  "location": {"longitude": 20.5, "latitude": 49.9},
  "data": {
    "windSpeed": 15.5
  }
}
\end{lstlisting}

\paragraph{simulation\_telemetry\_sensors\_wind\_direction}

\textbf{Routing Key:} \texttt{simulation.telemetry.sensors.wind\_direction}

\textbf{Event Model:} \texttt{EvWindDirectionSensor}

\textbf{Contract:}
\begin{lstlisting}[
  float,
  placement=H,
  caption={}
]
{
  "sensorId": 3,
  "location": {"longitude": 20.5, "latitude": 49.9},
  "data": {
    "windDirection": 270.0
  }
}
\end{lstlisting}

\textbf{Required:} \texttt{data.windDirection} (double): Kierunek w stopniach (0-360)

\paragraph{simulation\_telemetry\_sensors\_litter\_moisture}

\textbf{Routing Key:} \texttt{simulation.telemetry.sensors.litter\_moisture}

\textbf{Event Model:} \texttt{EvLitterMoistureSensor}

\textbf{Contract:}
\begin{lstlisting}[
  float,
  placement=H,
  caption={}
]
{
  "sensorId": 4,
  "location": {"longitude": 20.5, "latitude": 49.9},
  "data": {
    "litterMoisture": 45.0
  }
}
\end{lstlisting}

\paragraph{simulation\_telemetry\_sensors\_co2}

\textbf{Routing Key:} \texttt{simulation.telemetry.sensors.co2}

\textbf{Event Model:} \texttt{EvCO2Sensor}

\textbf{Contract:}
\begin{lstlisting}[
  float,
  placement=H,
  caption={}
] 
{
  "sensorId": 5,
  "location": {"longitude": 20.5, "latitude": 49.9},
  "data": {
    "co2Concentration": 420.0
  }
}
\end{lstlisting}

\paragraph{simulation\_telemetry\_sensors\_pm2\_5}

\textbf{Routing Key:} \texttt{simulation.telemetry.sensors.pm2\_5}

\textbf{Event Model:} \texttt{EvPM25ConcentrationSensor}

\textbf{Contract:}
\begin{lstlisting}[
  float,
  placement=H,
  caption={}
] 
{
  "sensorId": 6,
  "location": {"longitude": 20.5, "latitude": 49.9},
  "data": {
    "pm2_5Concentration": 25.5
  }
}
\end{lstlisting}

\paragraph{simulation\_telemetry\_sensors\_camera}

\textbf{Routing Key:} \texttt{simulation.telemetry.sensors.camera}

\textbf{Event Model:} \texttt{EvCamera}

\textbf{Contract:}
\begin{lstlisting}[
  float,
  placement=H,
  caption={}
] 
{
  "sensorId": 7,
  "location": {"longitude": 20.5, "latitude": 49.9},
  "data": {
    "smokeLevel": 15.0
  }
}
\end{lstlisting}

\textbf{Note:} Obecnie logowane, ale nie w pełni przetwarzane w backendzie.

\subsubsection{simulation\_recommendations}

\textbf{Routing Key:} \texttt{simulation.recommendations}

\textbf{Event Model:} \texttt{EvRecommendation}

\textbf{Contract:}
\begin{lstlisting}[
  float,
  placement=H,
  caption={}
] 
{
  "timestamp": "2026-01-14T17:30:53",
  "recommendedActions": [
    {
      "unitId": 1,
      "sectorId": 5,
      "actionType": "EXTINGUISH"
    }
  ]
}
\end{lstlisting}

\subsection{Kolejki wychodzące (TO Simulation/Support)}

\subsubsection{support\_data\_aggregated}

\textbf{Routing Key:} \texttt{support.aggregated.data}

\textbf{Payload:} \texttt{SimulationStateDto} (zoptymalizowany -- tylko zmienione encje)

\textbf{Contract:}
\begin{lstlisting}[
  float,
  placement=H,
  caption={}
] 
{
  "forestName": "Forest Name",
  "timestamp": "2026-01-14T17:30:53.000Z",
  "tick": 100,
  "sectors": [
    {
      "sectorId": 1,
      "state": {
        "fireLevel": 25.5,
        "burnLevel": 10.2,
        "extinguishLevel": 5.0,
        "fireState": "ACTIVE"
      },
      "contours": [[20.5, 49.9], [20.6, 49.9]],
      "assignedBrigades": [1, 2]
    }
  ],
  "fireBrigades": [ /* tylko zmienione brygady */ ],
  "foresterPatrols": [ /* tylko zmienione patrole */ ],
  "recommendedActions": [ /* tylko nowe rekomendacje */ ]
}
\end{lstlisting}

\textbf{Optimization:}
\begin{itemize}
    \item \textbf{Pierwsza wiadomość:} Pełny stan (wszystkie sektory, brygady, patrole)
    \item \textbf{Kolejne wiadomości:} Tylko zmienione encje
    \begin{itemize}
        \item Tylko sektory ze zmienionym fireLevel, burnLevel, extinguishLevel lub assignedBrigades
        \item Tylko brygady ze zmienioną lokalizacją, stanem lub sectorId
        \item Tylko patrole ze zmienioną lokalizacją, stanem lub sectorId
        \item Tylko nowe rekomendacje (nie wcześniej wysłane)
    \end{itemize}
\end{itemize}

\textbf{Performance Impact:} Redukcja transferu danych o 80-95\% w typowych scenariuszach.

\subsubsection{simulation\_control\_fire\_brigade\_actions}

\textbf{Routing Key:} \texttt{simulation.control.fire\_brigade\_actions}

\textbf{Payload:} \texttt{OrderFireBrigade}

\textbf{Contract:}
\begin{lstlisting}[
  float,
  placement=H,
  caption={}
] 
{
  "fireBrigadeId": 1,
  "action": "EXTINGUISH",
  "sectorId": null,
  "timestamp": "2026-01-14T17:30:53.000Z",
  "location": {
    "longitude": 20.5,
    "latitude": 49.9
  }
}
\end{lstlisting}

\textbf{Required Fields:}
\begin{itemize}
    \item \texttt{fireBrigadeId} (int): Brygada do rozkazania
    \item \texttt{action} (FireBrigadeAction): EXTINGUISH lub GO\_TO\_BASE
    \item \texttt{location} (Location): Docelowa lokalizacja (null dla GO\_TO\_BASE)
\end{itemize}

\subsubsection{simulation\_control\_forester\_actions}

\textbf{Routing Key:} \texttt{simulation.control.forester\_actions}

\textbf{Payload:} \texttt{OrderForesterPatrol}

\textbf{Contract:}
\begin{lstlisting}[
  float,
  placement=H,
  caption={}
] 
{
  "foresterPatrolId": 1,
  "action": "PATROL",
  "timestamp": "2026-01-14T17:30:53.000Z",
  "location": {
    "longitude": 20.5,
    "latitude": 49.9
  }
}
\end{lstlisting}

\subsubsection{support.recommendations}

\textbf{Routing Key:} \texttt{support.recommendations}

\textbf{Direction:} Publisher (wysyła do backend/frontend)

\textbf{Contract:}
\begin{lstlisting}[
  float,
  placement=H,
  caption={}
] 
{
  "timestamp": 1234567890.0,
  "recommendedActions": [
    {
      "unitId": 1,
      "sectorId": 5
    }
  ],
  "priority": "HIGH"
}
\end{lstlisting}

\textbf{Required Fields:}
\begin{itemize}
    \item \texttt{timestamp} (float): Unix timestamp w sekundach
    \item \texttt{recommendedActions} (List): Lista akcji (musi być niepusta)
    \item \texttt{priority} (str): "LOW", "MEDIUM" lub "HIGH"
\end{itemize}

\section{Formaty JSON}

\subsection{Configuration}

Format konfiguracji lasu używany do inicjalizacji symulacji.

\begin{lstlisting}[
  float,
  placement=H,
  caption={}
] 
{
  "forestName": "Test Forest",
  "rows": 10,
  "columns": 10,
  "location": [
    {"latitude": 49.9, "longitude": 20.5},
    {"latitude": 49.9, "longitude": 20.6},
    {"latitude": 49.8, "longitude": 20.6},
    {"latitude": 49.8, "longitude": 20.5}
  ],
  "sectors": [
    {
      "sectorId": 1,
      "row": 0,
      "column": 0,
      "sectorType": "FOREST",
      "initialState": {
        "fireLevel": 0,
        "burnLevel": 0,
        "extinguishLevel": 0,
        "fireState": "INACTIVE",
        "temperature": 20.0,
        "airHumidity": 60.0,
        "windSpeed": 10.0,
        "windDirection": 270.0,
        "litterMoisture": 50.0,
        "co2Concentration": 400.0,
        "pm2_5Concentration": 20.0
      },
      "contours": [
        [20.5, 49.9],
        [20.51, 49.9],
        [20.51, 49.89],
        [20.5, 49.89]
      ],
      "assignedBrigades": []
    }
  ],
  "fireBrigades": [
    {
      "fireBrigadeId": 1,
      "currentLocation": {
        "longitude": 20.5,
        "latitude": 49.9
      },
      "baseLocation": {
        "longitude": 20.5,
        "latitude": 49.9
      },
      "state": "IDLE"
    }
  ],
  "foresterPatrols": [
    {
      "foresterPatrolId": 1,
      "currentLocation": {
        "longitude": 20.6,
        "latitude": 49.8
      },
      "baseLocation": {
        "longitude": 20.6,
        "latitude": 49.8
      },
      "state": "PATROLLING"
    }
  ],
  "sensors": [
    {
      "sensorId": 1,
      "sensorType": "TEMP_AND_AIR_HUMIDITY",
      "location": {
        "longitude": 20.5,
        "latitude": 49.9
      }
    }
  ],
  "cameras": [
    {
      "cameraId": 1,
      "location": {
        "longitude": 20.5,
        "latitude": 49.9
      }
    }
  ]
}
\end{lstlisting}

\subsection{SimulationStateDto}

Format stanu symulacji wysyłany przez SSE.

\begin{lstlisting}[
  float,
  placement=H,
  caption={}
] 
{
  "forestName": "Test Forest",
  "timestamp": "2026-01-14T17:30:53.000Z",
  "tick": 100,
  "sectors": [
    {
      "sectorId": 1,
      "state": {
        "fireLevel": 25.5,
        "burnLevel": 10.2,
        "extinguishLevel": 5.0,
        "fireState": "ACTIVE"
      },
      "contours": [[20.5, 49.9], [20.6, 49.9]],
      "assignedBrigades": [1, 2]
    }
  ],
  "fireBrigades": [
    {
      "fireBrigadeId": 1,
      "location": {
        "longitude": 20.5,
        "latitude": 49.9
      },
      "state": "EXTINGUISHING",
      "sectorId": 5
    }
  ],
  "foresterPatrols": [
    {
      "foresterPatrolId": 1,
      "location": {
        "longitude": 20.6,
        "latitude": 49.8
      },
      "state": "PATROLLING",
      "sectorId": 3
    }
  ],
  "recommendedActions": [
    {
      "unitId": 1,
      "sectorId": 5
    }
  ]
}
\end{lstlisting}

\section{Obsługa błędów}

\subsection{Resource Not Found Errors}

Gdy zasób (sektor, brygada, patrol) nie zostanie znaleziony, błąd jest logowany z:
\begin{itemize}
    \item \textbf{Level:} ERROR
    \item \textbf{Format wiadomości:} \texttt{"Resource not found: \{ResourceType\} \{id\} does not exist when \{action\} at location \{location\}"}
    \item \textbf{Przykład:} \texttt{"Resource not found: Sector 5 does not exist when updating wind speed sensor at location Location[longitude=20.5, latitude=49.9]"}
\end{itemize}

Te błędy są logowane, ale nie powodują crashowania serwisu. Wiadomość jest pomijana i przetwarzanie kontynuowane.

\subsection{Sector Resolution}

Gdy \texttt{sectorId} nie jest podane w przychodzących wiadomościach:
\begin{enumerate}
    \item Sektor jest rozwiązywany z lokalizacji używając \texttt{SectorIdResolver}
    \item Jeśli rozwiązanie się nie powiedzie, błąd jest logowany i wiadomość pomijana
    \item Jeśli sektory nie mają jeszcze konturów, wiadomość jest odroczona do czasu dostępności konturów
\end{enumerate}

\section{Optymalizacje wydajności}

\subsection{Delta Updates}

Tylko zmienione encje są wysyłane do serwisu wsparcia, co redukuje ruch sieciowy o 80-95\%.

\subsection{Synchronizacja}

Wszystkie aktualizacje stanu używają bloków \texttt{synchronized(currentState)}, co zapobiega warunkom wyścigu. Rozważ użycie \texttt{ConcurrentHashMap} dla przyszłych ulepszeń.

\subsection{Message Filtering}

Wiadomości z czujników są pomijane, jeśli sektory nie mają jeszcze konturów, co zapobiega błędom podczas fazy inicjalizacji.

\section{Wersjonowanie i kompatybilność wsteczna}

\subsection{Obecny stan wersjonowania}

\textbf{Aktualna wersja systemu:} 1.0.0 (pre-release)

\textbf{Wersje API:}
\begin{itemize}
    \item \textbf{FFBackend:} Brak formalnego wersjonowania (bez prefiksu \texttt{/api/v1/})
    \item \textbf{FFConf:} Wersjonowane (\texttt{/api/v1/nodes})
    \item \textbf{RabbitMQ:} Brak pola \texttt{version} w wiadomościach
\end{itemize}

\subsection{Planowane wersjonowanie}

\subsubsection{REST API}

\textbf{Format:} \texttt{/api/v\{major\}/\{resource\}}

\textbf{Przykład:}
\begin{lstlisting}[
  float,
  placement=H,
  caption={}
]GET /api/v1/simulation/status
POST /api/v1/simulation/run-simulation
\end{lstlisting}

\textbf{Migracja:} Legacy endpointy (bez \texttt{/api/v1/}) będą działać jako alias do \texttt{v1} przez 6 miesięcy.

\subsubsection{RabbitMQ Messages}

\textbf{Format:} Pole \texttt{version} w każdej wiadomości

\textbf{Przykład:}
\begin{lstlisting}[
  float,
  placement=H,
  caption={}
] 
{
  "version": "1.0",
  "timestamp": "2026-01-14T17:30:53.000Z",
  "data": {
    "sectorId": 1,
    "fireLevel": 25.5
  }
}
\end{lstlisting}

\textbf{Alternatywa:} Routing key z wersją (np. \texttt{simulation.telemetry.v1.sectors})

\subsection{Breaking vs Non-Breaking Changes}

\subsubsection{Breaking Changes}

\textbf{Definicja:} Zmiana wymagająca aktualizacji kodu klienta.

\textbf{Przykłady:}
\begin{itemize}
    \item Usunięcie wymaganego pola (np. \texttt{fireBrigadeId})
    \item Zmiana typu pola (np. \texttt{sectorId: int} → \texttt{string})
    \item Zmiana nazwy pola (np. \texttt{fireLevel} → \texttt{intensity})
    \item Zmiana semantyki (np. \texttt{fireLevel} w skali 0-100 → 0-1)
    \item Zmiana routing key RabbitMQ
    \item Usunięcie endpointu API
\end{itemize}

\textbf{Procedura wprowadzenia:}
\begin{enumerate}
    \item Oznacz starą wersję jako \texttt{@Deprecated} (minimum 3 miesiące przed usunięciem)
    \item Dodaj nową wersję API równolegle (np. \texttt{/api/v2/})
    \item Dokumentuj zmianę w \texttt{CHANGELOG.md} i \texttt{MIGRATION\_GUIDE.md}
    \item Dodaj response header \texttt{X-Deprecated: true} do starych endpointów
    \item Loguj WARNING gdy stary endpoint jest używany
    \item Wyślij komunikat do użytkowników (email, dokumentacja)
    \item Usuń starą wersję po okresie deprecation
\end{enumerate}

\subsubsection{Non-Breaking Changes}

\textbf{Definicja:} Zmiana niewymagająca aktualizacji kodu klienta.

\textbf{Przykłady:}
\begin{itemize}
    \item Dodanie opcjonalnego pola (np. \texttt{extinguishProgress})
    \item Dodanie nowego endpointu (np. \texttt{/simulation/pause})
    \item Dodanie nowej kolejki RabbitMQ
    \item Rozszerzenie enum (np. \texttt{FireState} + \texttt{SMOLDERING})
    \item Zwiększenie limitu pola (np. \texttt{max fireLevel: 100} → \texttt{200})
\end{itemize}

\textbf{Procedura wprowadzenia:}
\begin{enumerate}
    \item Dodaj nowe pole/endpoint/kolejkę
    \item Dokumentuj w \texttt{CHANGELOG.md}
    \item Zachowaj backward compatibility (stare klienty działają bez zmian)
    \item Dodaj testy regresji
\end{enumerate}

\subsection{Deprecation Policy}

\subsubsection{Zasady}

\begin{itemize}
    \item \textbf{Minimalny okres wsparcia:} 3 miesiące od oznaczenia jako deprecated
    \item \textbf{Komunikacja:}
    \begin{itemize}
        \item Response header \texttt{X-Deprecated: true}
        \item Response header \texttt{X-Deprecated-Sunset: 2026-04-01}
        \item Warning w logach serwera
        \item Notice w dokumentacji API
    \end{itemize}
    \item \textbf{Dokumentacja:} Oznaczenie w API docs jako \texttt{[DEPRECATED -- use /api/v2/endpoint instead]}
\end{itemize}

\subsubsection{Przykład Deprecation Notice}

\textbf{HTTP Response:}
\begin{lstlisting}[
  float,
  placement=H,
  caption={}
]HTTP/1.1 200 OK
X-Deprecated: true
X-Deprecated-Sunset: 2026-04-01
X-Deprecated-Replacement: /api/v2/simulation/run-simulation
Content-Type: application/json

{
  "message": "This endpoint is deprecated and will be removed on 2026-04-01",
  "replacement": "/api/v2/simulation/run-simulation",
  "data": { ... }
}
\end{lstlisting}

\subsection{Migration Guide}

W przypadku breaking changes, dokumentacja powinna zawierać:

\subsubsection{Format Migration Guide}

\begin{lstlisting}[
  float,
  placement=H,
  caption={}
]# Migration Guide: v1 -> v2

## Breaking Changes

### 1. sectorId type changed: int -> string

**Before (v1):**
{
  "sectorId": 1,
  "fireLevel": 25.5
}

**After (v2):**
{
  "sectorId": "sector-001",
  "fireLevel": 25.5
}

**Migration:**
Convert sectorId to string with "sector-" prefix:
  v2_sectorId = "sector-" + str(v1_sectorId).zfill(3)

**Timeline:**
- v1 deprecated: 2026-02-01
- v1 removed: 2026-05-01
\end{lstlisting}

\subsection{Versioning dla różnych komponentów}

\subsubsection{REST API Versioning}

\textbf{Strategia:} URL Path Versioning

\textbf{Przykład:}
\begin{itemize}
    \item \texttt{/api/v1/simulation/run-simulation} (current)
    \item \texttt{/api/v2/simulation/run-simulation} (future)
\end{itemize}

\textbf{Alternatywy odrzucone:}
\begin{itemize}
    \item \textit{Header Versioning} (\texttt{Accept: application/vnd.fire-sim.v1+json}) -- zbyt skomplikowane dla małego projektu
    \item \textit{Query Parameter} (\texttt{?version=1}) -- trudniejsze do cachowania
\end{itemize}

\subsubsection{RabbitMQ Message Versioning}

\textbf{Strategia:} Message Field Versioning

\textbf{Przykład:}
\begin{lstlisting}[
  float,
  placement=H,
  caption={}
] {
  "version": "1.0",
  "data": { ... }
}
\end{lstlisting}

\textbf{Routing Key:} Bez wersji (stabilny: \texttt{simulation.telemetry.sectors})

\textbf{Uzasadnienie:} Routing key z wersją (\texttt{simulation.telemetry.v1.sectors}) wymaga rekonfiguracji exchanges przy każdej zmianie.

\subsubsection{JSON Schema Versioning}

\textbf{Strategia:} Schema ID w konfiguracji

\textbf{Przykład:}
\begin{lstlisting}[
  float,
  placement=H,
  caption={}
] {
  "$schema": "https://fire-sim.org/schemas/configuration/v1.json",
  "forestName": "Test Forest",
  ...
}
\end{lstlisting}

\subsection{Changelog Format}

\textbf{Standard:} Keep a Changelog + Semantic Versioning

\textbf{Przykład:}
\begin{lstlisting}[
  float,
  placement=H,
  caption={}
]# Changelog

## [1.1.0] - 2026-02-01

### Added
- FFBackend: POST /simulation/pause endpoint
- RabbitMQ: alerts.notifications queue
- FFSup: LLM coordination mode

### Changed
- FFSim: fireLevel calculation formula (non-breaking)
- FFBackend: Improved delta update performance

### Deprecated
- FFBackend: /orderFireBrigade (use /api/v1/commands/fire-brigade)
- RabbitMQ: simulation_telemetry_agents_fire_brigade queue name
  (use simulation.telemetry.agents.fire_brigade)

### Removed
- None

### Fixed
- FFBackend: SSE connection timeout after 60s
- FFSim: Race condition in sector state updates

### Security
- FFConf: Added input validation for node creation

## [1.0.0] - 2026-01-15

### Added
- Initial release
- All core modules (FFBackend, FFSim, FFSup, FFVis, FFConf)
- RabbitMQ integration
- MCTS decision engine
\end{lstlisting}

\section{Error Handling \& Retry Policies}

\subsection{HTTP Error Handling}

\subsubsection{Standardowy format błędu}

\begin{lstlisting}[
  float,
  placement=H,
  caption={}
] {
  "timestamp": "2026-01-14T17:30:53.000Z",
  "status": 503,
  "error": "Service Unavailable",
  "message": "Simulation service unreachable",
  "path": "/simulation/orderFireBrigade",
  "traceId": "abc123xyz" // opcjonalnie dla debugowania
}
\end{lstlisting}

\subsubsection{Kody błędów HTTP}

\begin{longtable}{|l|l|p{6cm}|}
\hline
\textbf{Kod} & \textbf{Nazwa} & \textbf{Kiedy używany} \\
\hline
400 & Bad Request & Nieprawidłowe dane wejściowe (np. \texttt{fireBrigadeId} < 0) \\
\hline
404 & Not Found & Zasób nie istnieje (np. brygada o ID=999) \\
\hline
409 & Conflict & Stan konfliktu (np. brygada już wykonuje akcję) \\
\hline
422 & Unprocessable Entity & Semantycznie nieprawidłowe (np. brygada poza mapą) \\
\hline
500 & Internal Server Error & Nieoczekiwany błąd serwera \\
\hline
503 & Service Unavailable & Zależny serwis niedostępny (np. FFSim offline) \\
\hline
504 & Gateway Timeout & Timeout komunikacji z serwisem (> 5s) \\
\hline
\end{longtable}

\subsection{Retry Policies}

\subsubsection{FFBackend → FFSim (HTTP)}

\textbf{Strategia:} Exponential Backoff + Fallback do RabbitMQ

\textbf{Parametry:}
\begin{itemize}
    \item \textbf{Próby:} 3 (initial + 2 retries)
    \item \textbf{Timeout:} 5s per request
    \item \textbf{Backoff:} 1s, 2s, 4s
    \item \textbf{Fallback:} Jeśli wszystkie próby nieudane → kolejka RabbitMQ
\end{itemize}

\textbf{Pseudokod:}
\begin{lstlisting}[language=Python]
def send_command_to_simulator(command):
    for attempt in [1, 2, 3]:
        try:
            response = http.post(
                url="http://simulator:5000/command",
                json=command,
                timeout=5
            )
            if response.status == 200:
                return response
        except Timeout:
            if attempt < 3:
                sleep(2 ** (attempt - 1))  # 1s, 2s, 4s
            else:
                # Fallback to RabbitMQ
                rabbitmq.publish(
                    queue="simulation.commands",
                    message=command
                )
                return {"status": "queued"}
\end{lstlisting}

\subsubsection{RabbitMQ Consumer Retry}

\textbf{Strategia:} Dead Letter Queue (DLQ) + Manual Requeue

\textbf{Parametry:}
\begin{itemize}
    \item \textbf{Próby:} 3 (przez \texttt{x-delivery-count} header)
    \item \textbf{Backoff:} 5s, 30s, 300s (przez TTL na DLQ)
    \item \textbf{DLQ:} \texttt{simulation.telemetry.dlq}
\end{itemize}

\textbf{Przepływ:}
\begin{enumerate}
    \item Wiadomość konsumowana z \texttt{simulation.telemetry.sectors}
    \item Jeśli przetwarzanie fail → \texttt{nack(requeue=false)}
    \item Wiadomość trafia do DLQ
    \item DLQ ma TTL (5s, 30s, 300s w zależności od \texttt{x-delivery-count})
    \item Po TTL wiadomość wraca do głównej kolejki
    \item Jeśli \texttt{x-delivery-count} > 3 → wiadomość trafia do \texttt{simulation.telemetry.failed}
\end{enumerate}

\subsubsection{Circuit Breaker Pattern}

\textbf{Cel:} Zapobieganie przeciążeniu FFSim przez FFBackend

\textbf{Stany:}
\begin{itemize}
    \item \textbf{CLOSED:} Normalny ruch (wszystkie requesty przechodzą)
    \item \textbf{OPEN:} Circuit otwarty (wszystkie requesty natychmiast failują → fallback do RabbitMQ)
    \item \textbf{HALF\_OPEN:} Testowanie (pojedyncze requesty testują czy serwis działa)
\end{itemize}

\textbf{Parametry:}
\begin{itemize}
    \item \textbf{Threshold:} 5 błędów z rzędu → OPEN
    \item \textbf{Timeout:} 30s w stanie OPEN → HALF\_OPEN
    \item \textbf{Success Threshold:} 2 udane requesty → CLOSED
\end{itemize}

\textbf{Implementacja:} Netflix Hystrix lub Resilience4j

\subsection{Idempotency}

\subsubsection{Problem}

Retry może spowodować duplikację komend (np. brygada otrzymuje 2x ten sam rozkaz).

\subsubsection{Rozwiązanie}

\textbf{Idempotency Key:} Każda komenda ma unikalny \texttt{commandId}

\begin{lstlisting}[
  float,
  placement=H,
  caption={}
] {
  "commandId": "uuid-abc-123",
  "fireBrigadeId": 1,
  "targetSectorId": 5,
  "command": "MOVE"
}
\end{lstlisting}

\textbf{Cache:} FFSim przechowuje \texttt{commandId} w cache (Redis lub in-memory) na 5 minut

\textbf{Logika:}
\begin{lstlisting}[language=Python]
def process_command(command):
    if cache.exists(command.commandId):
        return {"status": "already_processed"}
    
    cache.set(command.commandId, ttl=300)  # 5 min
    execute_command(command)
    return {"status": "ok"}
\end{lstlisting}

\subsection{Rate Limiting}

\subsubsection{FFBackend API}

\textbf{Limity:}
\begin{itemize}
    \item \textbf{Global:} 100 requests/second (wszystkie endpointy)
    \item \textbf{Per endpoint:} 
    \begin{itemize}
        \item \texttt{/simulation/run-simulation}: 1 request/5s (bo startuje nową symulację)
        \item \texttt{/simulation/orderFireBrigade}: 10 requests/second
    \end{itemize}
\end{itemize}

\textbf{Response gdy limit przekroczony:}
\begin{lstlisting}[
  float,
  placement=H,
  caption={}
]HTTP/1.1 429 Too Many Requests
Retry-After: 5
X-RateLimit-Limit: 10
X-RateLimit-Remaining: 0
X-RateLimit-Reset: 1642181053

{
  "error": "Rate limit exceeded",
  "message": "Max 10 requests per second"
}
\end{lstlisting}

\subsubsection{RabbitMQ Backpressure}

\textbf{Problem:} FFBackend nie nadąża z przetwarzaniem telemetrii → kolejka rośnie

\textbf{Rozwiązanie:}
\begin{itemize}
    \item \textbf{Queue Length Limit:} 10,000 wiadomości (po przekroczeniu stare są odrzucane)
    \item \textbf{TTL:} 60s (stare wiadomości są usuwane)
    \item \textbf{Prefetch:} 100 (consumer pobiera max 100 wiadomości naraz)
\end{itemize}

\subsection{Timeout Configuration}

\begin{longtable}{|l|l|p{5cm}|}
\hline
\textbf{Połączenie} & \textbf{Timeout} & \textbf{Uzasadnienie} \\
\hline
FFBackend → FFSim (HTTP) & 5s & Symulacja powinna odpowiedzieć szybko \\
\hline
FFVis → FFBackend (HTTP) & 10s & Użytkownik czeka max 10s \\
\hline
SSE Keep-Alive & 30s & Ping co 30s zapobiega disconnect \\
\hline
RabbitMQ Consumer & Brak & Asynchroniczne, brak timeout \\
\hline
\end{longtable}

\section{Konfiguracja RabbitMQ}

\subsection{Queue Configuration}

Wszystkie kolejki są:
\begin{itemize}
    \item \textbf{Durable:} false (zgodnie z konfiguracją symulatora)
    \item \textbf{Exchange Type:} topic
    \item \textbf{Exchange Name:} Konfigurowalne przez \texttt{RABBITMQ\_EXCHANGE} (domyślnie: \texttt{fire-simulation-exchange})
\end{itemize}

\subsection{Connection Settings}

\begin{itemize}
    \item \textbf{Host:} Konfigurowalne przez \texttt{RABBITMQ\_HOST} (domyślnie: \texttt{localhost})
    \item \textbf{Port:} Konfigurowalne przez \texttt{RABBITMQ\_PORT} (domyślnie: \texttt{5672})
    \item \textbf{Username:} Konfigurowalne przez \texttt{RABBITMQ\_USER}
    \item \textbf{Password:} Konfigurowalne przez \texttt{RABBITMQ\_PASS}
\end{itemize}

\section{Przyszłe ulepszenia}

\subsection{Planowane optymalizacje}

\begin{itemize}
    \item \textbf{Partial Sector Updates:} Wysyłanie tylko zmienionych pól (fireLevel, burnLevel, etc.) zamiast pełnego sektora
    \item \textbf{Compression:} Kompresja dużych wiadomości
    \item \textbf{Batching:} Grupowanie wielu małych aktualizacji w pojedynczą wiadomość
    \item \textbf{ConcurrentHashMap:} Zastąpienie bloków synchronized kolekcjami współbieżnymi
    \item \textbf{Backpressure:} Implementacja właściwej obsługi backpressure dla wysokich wskaźników wiadomości
\end{itemize}

\section{Tabele referencyjne}

\subsection{Stany agentów}

\begin{longtable}{|l|l|}
\hline
\textbf{Agent Type} & \textbf{States} \\
\hline
FireBrigade & AVAILABLE, TRAVELLING, EXTINGUISHING, IDLE \\
\hline
ForesterPatrol & AVAILABLE, TRAVELLING, PATROLLING \\
\hline
\end{longtable}

\subsection{Stany sektorów}

\begin{longtable}{|l|}
\hline
\textbf{FireState} \\
\hline
INACTIVE \\
\hline
ACTIVE \\
\hline
LOST \\
\hline
\end{longtable}

\subsection{Typy sektorów}

\begin{longtable}{|l|}
\hline
\textbf{SectorType} \\
\hline
FOREST \\
\hline
GRASS \\
\hline
WATER \\
\hline
BUILDING \\
\hline
\end{longtable}

\subsection{Typy czujników}

\begin{longtable}{|l|}
\hline
\textbf{SensorType} \\
\hline
TEMP\_AND\_AIR\_HUMIDITY \\
\hline
WIND\_SPEED \\
\hline
WIND\_DIRECTION \\
\hline
LITTER\_MOISTURE \\
\hline
CO2 \\
\hline
PM2\_5 \\
\hline
CAMERA \\
\hline
\end{longtable}

\end{document}
